\section{Virtualisierung}
Virtualisierung ist ein schon lange bestehendes Konzept. Bei Virtualisierung
wird auf einem physischen Computer durch Software eine andere, isolierte
Umgebung abstrahiert. Bekannt ist das Konzept beispielsweise durch Oracle
VirtualBox. Allerdings ist Virtualisierung als die Abstraktion von
Hardware Ressourcen durch Software auch der Grundstein für Cloudtechnologie.
Durch Virtualisierung wird Ressource Pooling überhaupt erst ermöglicht, sonst
könnten keine voneinander getrennten Umgebungen auf einem Server existieren.
\subsection{Wie nutzt eine Cloud Virtualisierung ?}
\subsubsection{Virtuelle Maschinen (Virtual Machine)}
Bei einer VM, kurz für Virtuelle Maschine, wird eine eigenständige,
isolierte Instanz auf einem physischen System abstrahiert. Dadurch kann man
beispielsweise auf einem Windows-PC ein Linux-Betriebssystem innerhalb einer
VM ausführen. In der Cloud wird dies für Ressource Pooling genutzt, also das
Aufteilen eines physischen Servers in mehrere kleinere, unabhängige Instanzen.
So werden Infrastructure as a Service (IaaS) und dynamische Skalierung
überhaupt erst ermöglicht.
\subsubsection{Virtuelles Netzwerk}
Durch ein virtuelles Netzwerk wird die Kommunikation zwischen VMs sowie
mit externen Systemen überhaupt erst ermöglicht. Durch Peering, also das
direkte Verbinden zweier virtueller Netzwerke miteinander, können Services
untereinander kommunizieren, ohne den Umweg über das öffentliche Internet
nehmen zu müssen. Durch VPNs, kurz für Virtual Private Network, wird zudem
eine verschlüsselte Direktverbindung zu On-Premises-Rechenzentren ermöglicht,
um sicher mit Cloud-Services zu kommunizieren.
\subsubsection{Speichervirtualisierung}
Bei der Speichervirtualisierung werden mehrere physische Speichermedien durch
Software zu einem einzigen logischen Speicherpool zusammengefasst und
abstrahiert. Das ermöglicht beispielsweise Services wie Google Drive, iCloud oder im Enterprise Kontext Azure Managed Disks.

\subsection{Wie Funktioniert Virtualisierung}
\subsubsection{Hypervisor}
Ein Hypervisor ist die Software die Arbeitspeicher, Rechenleistung und Speicher in Pools zusammenfasst. Dadurch können auf eine Physichen Maschine viele VMs erstellt werden. Der Hypervisor beötigt inige Komponenten auf Betriebssystemebene, zum Beispiel einen Speichermanager, Prozessplaner, Input/Output-Stack (I/O), Gerätetreiber, Sicherheitsmanager, Netzwerk-Stack. Der Hypervisor stellt den einzelnen virtuellen Maschinen die zugewiesenen Ressourcen zur Verfügung und verwaltet die Planung der VM-Ressourcen im Verhältnis zu den physischen Ressourcen. Die physische Hardware übernimmt nach wie vor die Ausführung, das heißt die CPU führt nach wie vor CPU-Befehle aus, wie sie beispielsweise von den VMs angefordert werden, während der Hypervisor die Planung übernimmt.

\begin{table}[h]
    \centering
    \begin{tabular}{|l|p{5cm}|p{5cm}|}
        \hline
        \textbf{Merkmal} & \textbf{Typ 1 (Bare-Metal)} & \textbf{Typ 2 (Gehostet)} \\
        \hline
        Ausführung & Direkt auf der Hardware & Auf einem Host-Betriebssystem \\
        \hline
        Host-OS & Keins, ersetzt das OS & Benötigt ein Host-OS \\
        \hline
        Ressourcenzuweisung & Direkt durch den Hypervisor & Über das Host-OS \\
        \hline
        Beispiele & VMware ESXi, KVM & Oracle VirtualBox, VMware Workstation \\
        \hline
    \end{tabular}
    \caption{Vergleich Hypervisor Typ 1 und Typ 2}
    \label{tab:hypervisor}
\end{table}